% chapters/06_scientific_computing.tex
% Scientific Computing and Computational Complexity

\section{Overview}

The theoretical results of Chapters 3--5 identify the FBSDE system, the
Fokker--Planck SPDE, and the $N$-particle system as the objects to be solved
numerically. This chapter addresses how to solve them \emph{at scale}: the
linear-algebra structures that arise, the iterative solvers best suited to them,
the parallel computing strategies that make large-$N$ and fine-mesh simulations
feasible, and the complexity of every algorithm, stated precisely in time, memory,
and parallel cost.

The chapter is organized as follows. Section~\ref{sec:sc-linalg} covers the
sparse linear-algebra structures arising from the finite-difference and particle
discretisations. Section~\ref{sec:sc-krylov} presents the Krylov subspace methods
(GMRES, CG, BiCG-STAB) used to solve the resulting linear systems. Section~\ref{sec:sc-precond}
discusses preconditioning strategies. Section~\ref{sec:sc-parallel} covers
parallel algorithms (OpenMP, MPI) and the scientific computing ecosystems
(PETSc, Trilinos) that implement them. Section~\ref{sec:sc-complexity} gives a
unified complexity table for every algorithm in Chapters 4--5.
Section~\ref{sec:sc-reproducibility} describes the reproducibility infrastructure
of this thesis.

\section{Sparse Linear-Algebra Structures}
\label{sec:sc-linalg}

\subsection{Finite-Difference Discretisation of the Fokker--Planck SPDE}

Discretising the Fokker--Planck SPDE (Chapter 4) on a uniform spatial grid
$x_j = j\Delta x$, $j = 0, \ldots, M$, yields at each time step a linear system
\begin{equation}
    \mathbf{A} \boldsymbol{\rho}^{n+1} = \mathbf{b}^n,
    \label{eq:sc-linear-system}
\end{equation}
where $\boldsymbol{\rho}^n = (\rho^n_0, \ldots, \rho^n_M)^\top$ is the density
vector and $\mathbf{A} \in \mathbb{R}^{(M+1)\times(M+1)}$ is a tridiagonal matrix:
\begin{equation}
    A_{jj} = 1 + \frac{\sigma^2 \Delta t}{\Delta x^2}, \quad
    A_{j,j\pm 1} = -\frac{\sigma^2 \Delta t}{2\Delta x^2} \pm
    \frac{b_j \Delta t}{2\Delta x},
    \label{eq:sc-tridiagonal}
\end{equation}
with the common-noise correction adding an $O(\sigma_0^2)$ perturbation to the
diagonal. The matrix $\mathbf{A}$ is:
\begin{itemize}
    \item \textit{Tridiagonal}: bandwidth 1, storage $O(M)$.
    \item \textit{Strictly diagonally dominant} under the CFL condition
    $\Delta t \le \Delta x^2 / \sigma^2$, guaranteeing $\mathbf{A}$ is invertible
    and the Thomas algorithm converges.
    \item \textit{M-matrix} (for upwind differencing), implying the discrete
    maximum principle holds: $\boldsymbol{\rho}^n \ge 0$ implies
    $\boldsymbol{\rho}^{n+1} \ge 0$.
\end{itemize}
\textbf{Thomas algorithm.} For tridiagonal systems, the Thomas algorithm
(a specialisation of Gaussian elimination) solves \eqref{eq:sc-linear-system}
in $O(M)$ operations with $O(M)$ memory, which is optimal.

\subsection{Two-Dimensional and Higher-Dimensional Cases}

For $d \ge 2$, the spatial grid has $M^d$ points and the system matrix
$\mathbf{A} \in \mathbb{R}^{M^d \times M^d}$ has bandwidth $M^{d-1}$.
\begin{itemize}
    \item \textit{Sparse storage}: $\mathbf{A}$ has $O(d M^d)$ nonzeros (each
    point connects to $2d$ neighbours). Compressed Sparse Row (CSR) format
    stores this in $O(d M^d)$ memory.
    \item \textit{Banded structure}: $\mathbf{A}$ is a $d$-point stencil matrix.
    For $d=2$, it is block tridiagonal with tridiagonal blocks (the ``five-point
    Laplacian'' structure).
    \item \textit{Direct solvers}: LU factorisation fills in and costs $O(M^{3d/2})$
    for $d=2$, prohibitive for large $M$. Iterative solvers are essential.
\end{itemize}

\subsection{HJB Equation Structure}

The backward HJB equation discretised on the same grid yields a system similar
to \eqref{eq:sc-linear-system} but with the drift reversed (the adjoint operator).
For the LQ-MFG, the HJB reduces to the Riccati ODE system of Appendix~A, which
is solved by a standard ODE integrator (RK45) rather than a linear system solver.
For general nonlinear HJB equations, Newton linearisation at each time step reduces
the problem to a sequence of sparse linear systems of the form
\eqref{eq:sc-linear-system}.

\section{Krylov Subspace Methods}
\label{sec:sc-krylov}

When direct solvers are too costly (i.e., $d \ge 2$ and $M$ large), Krylov
subspace methods solve $\mathbf{A}\mathbf{x} = \mathbf{b}$ iteratively by
building an approximation in the Krylov space
$\mathcal{K}_k(\mathbf{A}, \mathbf{b}) = \mathrm{span}\{\mathbf{b}, \mathbf{Ab},
\ldots, \mathbf{A}^{k-1}\mathbf{b}\}$.

\subsection{Conjugate Gradient (CG)}

When $\mathbf{A}$ is symmetric positive definite (SPD) --- which holds for the
diffusion-dominated Fokker--Planck operator after symmetrisation --- CG is optimal:

\begin{theorem}[CG convergence]
Let $\kappa = \lambda_{\max}(\mathbf{A}) / \lambda_{\min}(\mathbf{A})$ be the
condition number. After $k$ iterations, the CG error satisfies
\[
    \|\mathbf{x}_k - \mathbf{x}^*\|_{\mathbf{A}} \le 2\left(
    \frac{\sqrt{\kappa}-1}{\sqrt{\kappa}+1}\right)^k \|\mathbf{x}_0 -
    \mathbf{x}^*\|_{\mathbf{A}}.
\]
\end{theorem}

Each CG iteration costs one matrix-vector product ($O(M^d)$ for a sparse $d$-point
stencil), one dot product, and two vector updates: total $O(M^d)$ per iteration.
For the Fokker--Planck operator, $\kappa = O(\Delta x^{-2})$, so CG requires
$O(\Delta x^{-1})$ iterations to reduce the error by a fixed factor, giving total
cost $O(M^{d+1})$.

\subsection{GMRES}

When $\mathbf{A}$ is non-symmetric --- which is the generic case with drift terms
--- GMRES (Generalised Minimal RESidual) is the standard choice:
\[
    \mathbf{x}_k = \arg\min_{\mathbf{x} \in \mathbf{x}_0 +
    \mathcal{K}_k(\mathbf{A},\mathbf{r}_0)} \|\mathbf{A}\mathbf{x} - \mathbf{b}\|_2.
\]
GMRES is guaranteed to converge and minimises the residual over the Krylov space
at each step. Its cost is $O(k M^d)$ for $k$ iterations (storage of $k$ Krylov
vectors). Restarted GMRES($m$) limits memory to $O(m M^d)$ at the cost of
potentially slower convergence.

\subsection{BiCG-STAB}

BiConjugate Gradient Stabilised (BiCG-STAB) is a short-recurrence alternative to
GMRES for non-symmetric systems, with $O(M^d)$ memory per iteration (vs. $O(k M^d)$
for full GMRES). It is suitable when memory is the primary constraint.

\subsection{Choice of Solver for the MFG System}

\begin{table}[ht]
\centering
\caption{Recommended Krylov solver by operator type}
\label{tab:krylov-choice}
\resizebox{\textwidth}{!}{%
\begin{tabular}{llll}
\toprule
\textbf{Operator} & \textbf{Symmetry} & \textbf{Recommended solver} &
    \textbf{Notes} \\
\midrule
Pure diffusion (Laplacian) & SPD & CG & Optimal, $O(M^{d+1})$ \\
Diffusion + drift (FP) & Non-symmetric & GMRES or BiCG-STAB & Precondition
    with ILU \\
HJB (Newton step) & Non-symmetric & GMRES & Restart $m=50$ \\
Tridiagonal ($d=1$) & Either & Thomas algorithm & Direct, $O(M)$ \\
\bottomrule
\end{tabular}
}
\end{table}

\section{Preconditioning}
\label{sec:sc-precond}

The convergence rate of Krylov methods depends critically on the condition number
$\kappa(\mathbf{A})$. Preconditioning replaces $\mathbf{A}\mathbf{x}=\mathbf{b}$
with $\mathbf{M}^{-1}\mathbf{A}\mathbf{x} = \mathbf{M}^{-1}\mathbf{b}$ for a
preconditioner $\mathbf{M} \approx \mathbf{A}$ that is cheap to invert.

\subsection{Incomplete LU (ILU)}

ILU($k$) computes an approximate LU factorisation, retaining fill-in only up to
level $k$. For the sparse matrices arising here:
\begin{itemize}
    \item ILU(0): no fill-in; cost $O(M^d)$ to apply; typically reduces $\kappa$
    from $O(M^2)$ to $O(M)$ for diffusion-dominated problems.
    \item ILU(1): one level of fill; better conditioning at higher cost.
\end{itemize}

\subsection{Multigrid Preconditioning}

Geometric multigrid is theoretically optimal for elliptic operators: each
V-cycle reduces the error by a mesh-independent factor, and the total cost per
linear solve is $O(M^d)$. For the Fokker--Planck operator (which is an
advection-diffusion equation), algebraic multigrid (AMG, implemented in
\texttt{hypre} and \texttt{ML} within Trilinos) is robust and widely used.
AMG applied as a preconditioner for CG/GMRES gives $O(M^d)$ total cost,
independent of $\Delta x$.

\subsection{Block Diagonal Preconditioning for Coupled Systems}

The FBSDE system couples the forward (density) and backward (adjoint) equations.
A block diagonal preconditioner
\[
    \mathbf{M} = \begin{pmatrix} \mathbf{A}_{\mathrm{FP}} & 0 \\ 0 &
    \mathbf{A}_{\mathrm{HJB}} \end{pmatrix}
\]
decouples the two blocks and is effective when the coupling is weak (large $T$
or small Lipschitz constant $L$).

\section{Parallel Algorithms}
\label{sec:sc-parallel}

\subsection{Shared-Memory Parallelism: OpenMP}

For the particle system of Chapter 5, OpenMP parallelism is natural: each thread
updates a subset of particles independently (given fixed $\bar\mu_t^N, v_t^N$).

\begin{lstlisting}[language=C++, caption={OpenMP parallelisation of particle update.}]
#pragma omp parallel for schedule(static)
for (int i = 0; i < N; ++i) {
    double drift = kappa * (mu_bar - X[i]) - (A * X[i] + B * mu_bar) / lambda;
    X[i] += drift * dt + sigma * dW_idio[i] + sigma0 * dW_common;
}
// Reduction: compute new mu_bar and variance
double mu_bar_new = 0.0;
#pragma omp parallel for reduction(+:mu_bar_new)
for (int i = 0; i < N; ++i) mu_bar_new += X[i];
mu_bar_new /= N;
\end{lstlisting}

\textbf{Speedup.} For $P$ cores, the per-step cost reduces from $O(N)$ to
$O(N/P + \log P)$ (the $\log P$ term comes from the parallel reduction for
$\bar\mu_t^N$). Amdahl's law limits the speedup: if the serial fraction is $s$,
the maximum speedup is $1/s$ regardless of $P$.

\subsection{Distributed-Memory Parallelism: MPI}

For fine-mesh PDE solvers ($M^d$ large), MPI distributes the spatial domain across
processes. Each process owns a contiguous block of grid points; boundary communication
requires exchanging ``halo'' layers with neighbouring processes.

For the Fokker--Planck SPDE on $P$ MPI processes:
\begin{itemize}
    \item \textit{Subdomain size}: $M^d / P$ points per process.
    \item \textit{Communication}: $O(M^{d-1})$ boundary values per process per
    step (the halo layer).
    \item \textit{Parallel efficiency}: $(M^d/P) / (M^d/P + M^{d-1}) \to 1$
    as $M \to \infty$ for fixed $P$ (computation dominates communication).
\end{itemize}
The common noise $dW_t^0$ is a scalar broadcast ($O(\log P)$ with a tree
reduction) — negligible overhead.

\subsection{Scientific Computing Ecosystems}

\textbf{PETSc} (Portable, Extensible Toolkit for Scientific Computation) provides:
\begin{itemize}
    \item Sparse matrix formats (AIJ, BAIJ) with parallel assembly;
    \item Krylov solvers (CG, GMRES, BiCG-STAB) with unified interface;
    \item Preconditioners (ILU, AMG via \texttt{hypre});
    \item Time steppers (RK, IMEX) for the FBSDE system;
    \item Native MPI support with scalable parallel linear algebra.
\end{itemize}

\textbf{Trilinos} (Sandia National Laboratories) provides:
\begin{itemize}
    \item \texttt{Tpetra}: parallel sparse linear algebra;
    \item \texttt{Belos}: block and recycling Krylov solvers;
    \item \texttt{ML} / \texttt{MueLu}: algebraic multigrid preconditioners;
    \item \texttt{Zoltan}: dynamic load balancing for particle methods.
\end{itemize}

For the MFG system at research scale ($M \le 10^3$, $N \le 10^5$), PETSc with
GMRES + ILU(0) is sufficient. At production scale ($M \sim 10^4$, $N \sim 10^7$),
AMG preconditioning and MPI parallelism via PETSc or Trilinos are necessary.

\section{Complexity Analysis}
\label{sec:sc-complexity}

This section provides a unified complexity analysis for every algorithm introduced
in Chapters 4 and 5. We distinguish three cost measures: \textit{time} (sequential
floating-point operations), \textit{memory} (storage), and \textit{parallel time}
on $P$ processors (assuming perfect load balance).

\begin{table}[ht]
\centering
\caption{Algorithm complexity: $M$ spatial points, $N$ particles, $K$ time
steps, $d$ space dimension, $P$ processors, $k$ Krylov iterations.}
\label{tab:complexity}
\resizebox{\textwidth}{!}{%
\begin{tabular}{lllll}
\toprule
\textbf{Algorithm} & \textbf{Time (sequential)} & \textbf{Memory} &
    \textbf{Parallel time} & \textbf{Notes} \\
\midrule
Thomas algorithm ($d=1$) & $O(MK)$ & $O(M)$ & $O(MK/P)$ &
    Optimal for tridiagonal \\
GMRES (no precond.) & $O(kM^dK)$ & $O(kM^d)$ & $O(kM^dK/P + kK\log P)$ &
    $k$ Krylov its./step \\
CG + ILU(0) & $O(k_{\mathrm{ilu}} M^d K)$ & $O(M^d)$ & $O(k_{\mathrm{ilu}} M^d K/P)$ &
    $k_{\mathrm{ilu}} = O(M)$ for diffusion \\
CG + AMG & $O(M^d K)$ & $O(M^d)$ & $O(M^d K / P)$ &
    Optimal; $k_{\mathrm{amg}} = O(1)$ \\
\midrule
Euler--Maruyama (particles) & $O(NK)$ & $O(N)$ & $O(NK/P + K\log P)$ &
    LQ mean-field factor. \\
Particle (naive interact.) & $O(N^2K)$ & $O(N)$ & $O(N^2K/P)$ &
    General interaction \\
Particle (random batch $K_b$) & $O(NK_bK)$ & $O(N)$ & $O(NK_bK/P)$ &
    Approx.; batch size $K_b \ll N$ \\
Sinkhorn OT ($\varepsilon$, $d=1$) & $O(N\log N \cdot K)$ & $O(N)$ &
    $O(N\log N \cdot K/P)$ & Exact quantile formula \\
Sinkhorn OT ($\varepsilon$, $d\ge2$) & $O(N^2 K / \varepsilon^d)$ & $O(N^2)$ &
    $O(N^2 K / (\varepsilon^d P))$ & Regularised OT \\
\midrule
Full FBSDE (combined) & $O((M^d + N)K)$ & $O(M^d + N)$ &
    $O((M^d+N)K/P + K\log P)$ & With AMG + LQ factor. \\
Full FBSDE (naive) & $O(M^d K k + N^2 K)$ & $O(kM^d + N)$ & --- &
    Without precond./factor. \\
\midrule
LQ benchmark (Riccati) & $O(K_{\mathrm{ode}})$ & $O(1)$ & $O(1)$ &
    RK45 ODE; $K_{\mathrm{ode}} \sim 10^4$ \\
\bottomrule
\end{tabular}
}
\end{table}

\subsection{Dominant Cost and Bottlenecks}

For the typical parameter regime of Chapter 4 ($d=1$, $M=10^3$, $N=10^5$,
$K=10^3$):
\begin{itemize}
    \item \textit{PDE solve}: Thomas algorithm, $O(MK) = O(10^6)$ operations
    — negligible.
    \item \textit{Particle update}: $O(NK) = O(10^8)$ operations — dominant
    at $d=1$.
    \item \textit{Wasserstein computation}: $O(N\log N \cdot K) \approx O(10^9)$
    — the actual bottleneck for $d=1$.
\end{itemize}

For $d=2$ with $M=10^2$ per dimension ($10^4$ grid points) and $N=10^4$:
\begin{itemize}
    \item \textit{PDE solve with AMG}: $O(M^2 K) = O(10^9)$ — dominant.
    \item \textit{Particle update}: $O(NK) = O(10^7)$ — cheap.
    \item \textit{Sinkhorn OT}: $O(N^2 K / \varepsilon^2) \approx O(10^{12})$
    for $\varepsilon=0.1$ — prohibitive without approximation.
\end{itemize}
This confirms that Wasserstein distance computation is the scalability bottleneck
for $d \ge 2$, motivating the moment-embedding approximations discussed in
Chapter 8.

\section{Optimal Transport: Additional Structures}
\label{sec:sc-ot}

The Wasserstein distance computation is central both to the theoretical analysis
(Chapters 3--5) and to the numerical implementation. This section develops the
optimal transport structures beyond what was covered in Chapter 2.

\subsection{Benamou--Brenier Formulation}

Benamou and Brenier \citep{benamoubrenier2000} gave a dynamic characterisation
of the Wasserstein-2 distance:
\begin{equation}
    W_2(\mu_0, \mu_1)^2 = \min_{(\rho, v)} \int_0^1 \int_{\mathbb{R}^d}
    |v(t,x)|^2 \rho(t,x)\, dx\, dt,
    \label{eq:bb-formulation}
\end{equation}
where the minimum is over pairs $(\rho, v)$ satisfying the continuity equation
$\partial_t \rho + \nabla \cdot (\rho v) = 0$ with $\rho(0,\cdot) = \mu_0$
and $\rho(1,\cdot) = \mu_1$. This is a convex optimisation problem in $(\rho, v)$
and can be solved by operator-splitting algorithms (the ALG2 method of Benamou
and Brenier).

\textbf{Significance for MFGs.} The Benamou--Brenier formulation connects
the Wasserstein geometry to fluid dynamics: the optimal $v$ is the gradient of
a convex potential (Brenier's theorem, Chapter 2), and the continuity equation
is the Fokker--Planck equation without diffusion. The MFG Fokker--Planck equation
can therefore be viewed as a regularised version of the Benamou--Brenier dynamics,
with the $\sigma^2/2 \cdot \Delta\rho$ term providing entropic regularisation.

\subsection{JKO Scheme (Jordan--Kinderlehrer--Otto)}

The JKO scheme \citep{jordankinderlehrerotto1998} is a time-discretisation of
gradient flows in Wasserstein space. For a free energy functional
$\mathcal{F}(\rho) = \int (\rho \log \rho + V\rho)\, dx$, the gradient flow
$\partial_t \rho = -\mathrm{grad}_W \mathcal{F}(\rho)$ is approximated by the
implicit Euler scheme:
\begin{equation}
    \rho^{n+1} = \arg\min_{\rho \in \mathcal{P}_2} \left\{
    \frac{W_2(\rho, \rho^n)^2}{2\tau} + \mathcal{F}(\rho) \right\},
    \label{eq:jko-scheme}
\end{equation}
where $\tau > 0$ is the time step. The Fokker--Planck equation
$\partial_t \rho = \sigma^2/2 \cdot \Delta\rho + \nabla\cdot(\rho\nabla V)$
is the Wasserstein gradient flow of the free energy $\mathcal{F}$.

\textbf{Significance for numerical MFGs.} The JKO scheme provides an
alternative discretisation of the Fokker--Planck component of the MFG system
that automatically preserves positivity and mass conservation (since each step is
a Wasserstein projection). Boffi and Vanden-Eijnden (2023) used JKO-type schemes
for McKean--Vlasov equations; extending this to the common-noise case is an open
problem noted in Chapter 8.

\subsection{Otto Calculus and Gradient Flows}

Otto \citep{otto2001} showed that $(\mathcal{P}_2(\mathbb{R}^d), W_2)$ is
formally a Riemannian manifold (an infinite-dimensional one), with the tangent
space at $\mu$ identified with $L^2(\mu)$-gradients of scalar functions. The
\emph{Otto calculus} defines:
\begin{itemize}
    \item \textit{Riemannian metric:} $\langle v, w \rangle_\mu = \int v \cdot w\,
    d\mu$ for tangent vectors $v, w \in L^2(\mu)$.
    \item \textit{Gradient of $\mathcal{F}$:} $\mathrm{grad}_W \mathcal{F}(\mu)
    = -\nabla \cdot (\mu \nabla (\delta\mathcal{F}/\delta\mu))$ (the continuity
    equation direction).
    \item \textit{Geodesics:} straight lines in the Brenier-map parametrisation,
    $\rho_t = ((1-t)\mathrm{id} + t T)_\# \mu_0$ where $T = \nabla\varphi$ is
    the optimal transport map.
\end{itemize}
The Lions derivative $D_\mu U$ introduced in Chapter 2 is precisely the
``gradient'' of $U$ in the Otto calculus sense: $D_\mu U(t,x,\mu)(y)$ is the
derivative of $U$ along the geodesic direction at $y$.

\section{Reproducibility}
\label{sec:sc-reproducibility}

All numerical results in this thesis are reproducible. The following
infrastructure is provided:

\subsection{Code Repository}

The repository \texttt{https://github.com/NisongMonyimba/AppliedMathThesis}
contains:
\begin{itemize}
    \item \texttt{manuscript/}: all \LaTeX\ source files for this thesis;
    \item \texttt{simulations/}: Python simulation scripts for all numerical
    experiments:
    \begin{itemize}
        \item \texttt{euler\_milstein/}: Euler--Maruyama convergence benchmark
        (Chapter 4);
        \item \texttt{particle\_filter/}: particle approximation experiments
        (Chapter 5);
        \item \texttt{sobol\_sensitivity/}: sensitivity analysis scripts;
        \item \texttt{rl\_training/}: (archived) earlier RL experiments;
    \end{itemize}
    \item \texttt{lean4\_verification/}: Lean 4 source files for all formal
    proofs (Chapter 7);
    \item \texttt{.github/workflows/}: CI workflows that compile the \LaTeX\
    manuscript and run \texttt{lake build} for the Lean 4 proofs on every push.
\end{itemize}

\subsection{Numerical Experiment Reproducibility}

All simulation scripts:
\begin{itemize}
    \item set random seeds explicitly (\texttt{np.random.seed(42)}) for
    reproducibility;
    \item use common random numbers (CRN) for convergence rate estimation, so
    that stochastic error is isolated from discretisation error;
    \item write results to \texttt{.csv} files (e.g.,
    \texttt{convergence\_results.csv}) that are tracked in the repository;
    \item document all parameter choices in the script docstrings.
\end{itemize}
The Riccati sign corrections documented in Appendix A are reflected in all
four simulation scripts (\texttt{euler\_milstein}, \texttt{particle\_filter},
\texttt{sobol\_sensitivity}, \texttt{rl\_training}), with the corrected equation
and a comment indicating the fix.

\subsection{Lean 4 Proofs}

The Lean 4 proofs (Chapter 7) are built by the \texttt{lean4-check.yml}
CI workflow on every push to \texttt{main}. To rebuild locally:
\begin{lstlisting}[language=bash, caption={Lean 4 local build instructions.}]
cd lean4_verification/
lake update     # fetch Mathlib at pinned version
lake build      # compile all proofs
\end{lstlisting}
Known issues: Mathlib lemma names evolve; if \texttt{lake build} fails, check
the CI log for the current lemma name and update accordingly.

\subsection{Environment}

\begin{table}[ht]
\centering
\caption{Software environment for reproducibility}
\label{tab:environment}
\begin{tabular}{ll}
\toprule
\textbf{Component} & \textbf{Version / Details} \\
\midrule
Operating system & Ubuntu 24.04 (WSL2 on Windows 11) \\
Python & 3.11 (conda, \texttt{fenicsx} environment) \\
NumPy & 1.26 \\
SciPy & 1.12 (LSODA, Radau ODE solvers) \\
\LaTeX & TeX Live 2022, pdflatex \\
Lean 4 & v4.x (see \texttt{lean-toolchain} in repo) \\
Mathlib & pinned in \texttt{lakefile.lean} \\
\bottomrule
\end{tabular}
\end{table}

\section{Conclusion}

This chapter has developed the scientific computing foundation for the algorithms
of Chapters 4 and 5:
\begin{itemize}
    \item the sparse tridiagonal/block-sparse structure of the Fokker--Planck
    and HJB linear systems, and why direct solvers (Thomas, LU) are optimal
    for $d=1$ but Krylov methods are necessary for $d \ge 2$;
    \item GMRES, CG, and BiCG-STAB as the appropriate Krylov solvers for
    non-symmetric, symmetric, and memory-constrained settings respectively;
    \item ILU and AMG as preconditioners, with AMG giving $O(M^d)$ total cost;
    \item OpenMP and MPI parallelism with PETSc/Trilinos as the production
    ecosystems;
    \item a unified complexity table for every algorithm in the thesis;
    \item the Benamou--Brenier, JKO, and Otto calculus perspectives on optimal
    transport that deepen the connection between the Fokker--Planck dynamics
    and Wasserstein geometry;
    \item and the full reproducibility infrastructure of the repository.
\end{itemize}
Chapter 7 turns to the formal verification of the mathematical estimates
underlying these algorithms.
